\begin{table}[htbp]
\centering
\small
\caption{What the decision tree contributes. \textbf{No tree} ranks constructions
by their raw train-pool mean score. \textbf{Tree} is the estimator reported in
Table~\ref{tab:dt_consistency}, evaluated on the same variants it was fit on.
\textbf{LOVO} refits the tree with one construction held out and predicts that
construction, the only setting in which the tree ranks a variant it has not
scored. All three are Spearman $\rho$ against the held-out test-pool means. The
tree matches the no-tree floor, so the headline $\rho$ measures the
reproducibility of the proxy scores; the tree's contribution is compression of
the score table into readable rules, and its generalization to unscored
constructions is given by the LOVO column.}
\label{tab:tree_ablation}
\vspace{4pt}
\begin{tabular}{lrccc}
\toprule
Dataset & $n$ & No tree & Tree & LOVO \\
\midrule
\multicolumn{5}{c}{\textit{Node classification}} \\
\midrule
History & 8 & 0.952 & 0.945 & 0.452 \\
Sports & 26 & 0.954 & 0.939 & 0.881 \\
ArXiv & 16 & 0.959 & 0.958 & 0.868 \\
Electronics & 26 & 0.960 & 0.957 & 0.930 \\
Toys & 26 & 0.869 & 0.873 & 0.878 \\
\textbf{Mean} & & \textbf{0.939} & \textbf{0.934} & \textbf{0.802} \\
\midrule
\multicolumn{5}{c}{\textit{GraphRAG retrieval}} \\
\midrule
History & 8 & 0.929 & 0.976 & 0.786 \\
Sports & 26 & 0.802 & 0.814 & 0.764 \\
ArXiv & 16 & 0.853 & 0.840 & 0.544 \\
Electronics & 26 & 0.774 & 0.805 & 0.625 \\
Toys & 26 & 0.934 & 0.941 & 0.743 \\
\textbf{Mean} & & \textbf{0.858} & \textbf{0.875} & \textbf{0.692} \\
\bottomrule
\end{tabular}
\end{table}
